\documentclass[10pt,a4paper]{article}

\usepackage[T1]{fontenc}
\usepackage[utf8]{inputenc}
\usepackage[ngerman]{babel}
\usepackage[a4paper,margin=14mm]{geometry}
\usepackage{lmodern}
\usepackage[table]{xcolor}
\usepackage{tabularx}
\usepackage{array}
\usepackage{amsmath}
\usepackage{microtype}
\usepackage[most]{tcolorbox}
\usepackage{tikz}
\usetikzlibrary{arrows.meta,calc}

\pagestyle{empty}
\setlength{\parindent}{0pt}
\setlength{\parskip}{0pt}
\renewcommand{\familydefault}{\sfdefault}
\renewcommand{\arraystretch}{1.08}

\definecolor{Navy}{HTML}{17324D}
\definecolor{Blue}{HTML}{3B6EA5}
\definecolor{Sky}{HTML}{EAF3FA}
\definecolor{Mint}{HTML}{E8F5EF}
\definecolor{Gold}{HTML}{F3B33D}
\definecolor{Ink}{HTML}{1E2935}
\definecolor{Muted}{HTML}{617080}
\definecolor{Line}{HTML}{D8E1E8}
\colorlet{Force}{red!80!black}

\color{Ink}
\newcommand{\sectiontitle}[1]{%
  \vspace{1.0mm}{\color{Blue}\bfseries\large #1}\par
  \vspace{0.15mm}{\color{Blue}\rule{\linewidth}{0.7pt}}\par
  \vspace{0.6mm}}
\newcommand{\unit}[1]{\,\mathrm{#1}}
\newcommand{\unitfrac}[2]{\,\frac{\mathrm{#1}}{\mathrm{#2}}}
\newcommand{\vB}{v_{\mathrm B}}
\newcommand{\FZ}{F_{\mathrm Z}}
% Formelkasten: Formel groß, darunter die Größen mit Einheit.
\newcommand{\formulabox}[2]{%
  \begin{tcolorbox}[colback=Sky,colframe=Line,arc=1.2mm,boxrule=.6pt,left=2mm,right=2mm,top=.8mm,bottom=.8mm]
  {\centering\fontsize{15}{18}\selectfont #1\par}
  \vspace{0.3mm}{\centering\footnotesize #2\par}
  \end{tcolorbox}}
% Beispiel aus der Simulation des Lernmoduls.
\newcommand{\example}[1]{%
  \vspace{0.8mm}{\small\textcolor{Blue}{\bfseries Beispiel aus der Simulation:}\enspace #1\par}\vspace{0.8mm}}
% Zwei Je-desto-Kästen nebeneinander.
\newcommand{\twofacts}[2]{%
  \begin{tabularx}{\linewidth}{@{}X@{\hspace{1.5mm}}X@{}}
  \colorbox{Mint}{\parbox{\dimexpr\linewidth-2\fboxsep\relax}{\centering\small\vspace{0.4mm}#1\vspace{0.4mm}}} &
  \colorbox{Mint}{\parbox{\dimexpr\linewidth-2\fboxsep\relax}{\centering\small\vspace{0.4mm}#2\vspace{0.4mm}}}
  \end{tabularx}}

\begin{document}

\colorbox{Navy}{%
  \parbox{\dimexpr\linewidth-2\fboxsep\relax}{%
    \vspace{1.3mm}\color{white}
    \begin{tabularx}{\linewidth}{@{}Xr@{}}
      {\small\bfseries PHYSIK · KLASSE 11} &
      \colorbox{Gold}{\color{Navy}\small\bfseries\strut LÖSUNGEN}
    \end{tabularx}
    \vspace{0.9mm}
    {\fontsize{16}{19}\selectfont\bfseries Aufgabe 3 -- Winkelgeschwindigkeit und Ursache der Kreisbewegung}\par
    \vspace{0.3mm}{\large Sicherungsblatt}\vspace{1.1mm}
  }%
}

\vspace{0.9mm}
\colorbox{Sky}{\parbox{\dimexpr\linewidth-2\fboxsep\relax}{%
  \vspace{0.55mm}\textcolor{Blue}{\bfseries Ziel:} Eine Drehbewegung mit $\omega$, $\vB$, $f$ und $T$ beschreiben und erklären, warum ein Körper auf der Kreisbahn bleibt.\vspace{0.55mm}}}

% ---------------------------------------------------------------------------
\sectiontitle{1 · Winkelgeschwindigkeit}
\begin{minipage}[c]{0.27\linewidth}\centering
\begin{tikzpicture}[x=1mm,y=1mm]
  % Drehung von oben im Uhrzeigersinn wie in der Simulation.
  \fill[Sky] (0,0) -- (90:16) arc[start angle=90,end angle=0,radius=16] -- cycle;
  \draw[Muted,line width=.6pt] (0,0) circle (16);
  \draw[Muted,densely dashed,line width=.6pt] (0,0) -- (90:16);
  \draw[Navy,line width=.9pt] (0,0) -- (0:16);
  \draw[-{Latex[length=1.8mm]},Blue,line width=.9pt] (90:8) arc[start angle=90,end angle=4,radius=8];
  \node[Blue,font=\small\bfseries] at (45:4.6) {$\Delta\varphi$};
  \fill[Navy] (0,0) circle (.8);
  \node[font=\scriptsize,anchor=north] at (0,-1.2) {Zentrum};
  \fill[Muted!45] (90:16) circle (1.8);
  \fill[Navy] (0:16) circle (1.8);
  \node[font=\scriptsize,Muted,anchor=south] at (90:18) {Start};
  \node[font=\scriptsize,anchor=west] at (0:18.2) {nach $\Delta t$};
\end{tikzpicture}
\end{minipage}\hfill
\begin{minipage}[c]{0.70\linewidth}
\formulabox{$\omega=\dfrac{\Delta\varphi}{\Delta t}$}%
  {$\Delta\varphi$: Drehwinkel im Bogenmaß (rad) \quad $\Delta t$: Zeitintervall (s) \quad $\omega$: Winkelgeschwindigkeit $\left(\unitfrac{rad}{s}\right)$}
\example{$\Delta\varphi=\dfrac{\pi}{2}$, $\Delta t=1{,}0\unit{s}$ \quad$\Rightarrow$\quad $\omega=\dfrac{\pi/2}{1{,}0\unit{s}}\approx1{,}6\unitfrac{rad}{s}$}
\twofacts{Gleicher Winkel $\Delta\varphi$ in \textbf{kürzerer} Zeit\\$\to$ \textbf{größeres} $\omega$}%
         {\textbf{Kleinerer} Winkel $\Delta\varphi$ in gleicher Zeit\\$\to$ \textbf{kleineres} $\omega$}
\end{minipage}

% ---------------------------------------------------------------------------
\sectiontitle{2 · Bahngeschwindigkeit}
\begin{minipage}[c]{0.27\linewidth}\centering
\begin{tikzpicture}[x=1mm,y=1mm]
  % Gleiches omega: Beide Körper liegen auf derselben Verbindungslinie.
  \draw[Muted,line width=.6pt] (0,0) circle (8);
  \draw[Muted,line width=.6pt] (0,0) circle (16);
  \draw[Navy,line width=.7pt] (0,0) -- (0,16);
  \draw[-{Latex[length=2mm]},Blue,line width=1.1pt] (0,8) -- (9,8);
  \draw[-{Latex[length=2mm]},Blue,line width=1.1pt] (0,16) -- (18,16);
  \fill[Navy] (0,8) circle (1.5);
  \fill[Navy] (0,16) circle (1.5);
  \fill[Navy] (0,0) circle (.8);
  \node[Blue,font=\scriptsize,anchor=west] at (9.2,8) {$\vB$};
  \node[Blue,font=\scriptsize,anchor=west] at (18.2,16) {$\vB$};
  \node[font=\scriptsize,anchor=north] at (0,-1.2) {Zentrum};
  \node[font=\scriptsize,Muted,anchor=north] at (0,-17.2) {gleiches $\omega$};
\end{tikzpicture}
\end{minipage}\hfill
\begin{minipage}[c]{0.70\linewidth}
\formulabox{$\vB=\omega\cdot r$}%
  {$\omega$: Winkelgeschwindigkeit $\left(\unitfrac{rad}{s}\right)$ \quad $r$: Radius (m) \quad $\vB$: Bahngeschwindigkeit $\left(\unitfrac{m}{s}\right)$}
\example{$r=2{,}0\unit{m}$, $\omega=2{,}0\unitfrac{rad}{s}$ \quad$\Rightarrow$\quad $\vB=2{,}0\unitfrac{rad}{s}\cdot2{,}0\unit{m}=4{,}0\unitfrac{m}{s}$}
\twofacts{Gleicher Radius $r$, \textbf{größeres} $\omega$\\$\to$ \textbf{größeres} $\vB$}%
         {Gleiches $\omega$, \textbf{größerer} Radius $r$\\$\to$ \textbf{größeres} $\vB$}
\end{minipage}

% ---------------------------------------------------------------------------
\sectiontitle{3 · Frequenz und Umlaufdauer}
\begin{minipage}[c]{0.40\linewidth}
\formulabox{$f=\dfrac{1}{T}$}%
  {$T$: Umlaufdauer (s)\\[0.3mm] $f$: Frequenz in Hertz $\left(1\unit{Hz}=\tfrac{1}{\mathrm{s}}\right)$}
\end{minipage}\hfill
\begin{minipage}[c]{0.57\linewidth}
\example{$T=2{,}0\unit{s}$ \quad$\Rightarrow$\quad $f=\dfrac{1}{2{,}0\unit{s}}=0{,}50\unit{Hz}$}
\colorbox{Mint}{\parbox{\dimexpr\linewidth-2\fboxsep\relax}{\centering\small\vspace{0.4mm}$f$ und $T$ sind \textbf{Kehrwerte}: \textbf{kurze} Umlaufdauer $\to$ viele Umläufe pro Sekunde $\to$ \textbf{große} Frequenz\vspace{0.4mm}}}
\end{minipage}
\par\vspace{1.5mm}

% ---------------------------------------------------------------------------
\sectiontitle{4 · Ursache der Kreisbewegung}
\begin{minipage}[c]{0.24\linewidth}\centering
\begin{tikzpicture}[x=1mm,y=1mm]
  \coordinate (M) at (0,0);
  \coordinate (K) at (45:15);
  \draw[Muted,line width=.6pt] (M) circle (15);
  \fill[Navy] (M) circle (.8);
  \node[font=\scriptsize,anchor=north] at (0,-1.2) {Mittelpunkt};
  \draw[-{Latex[length=2mm]},Force,line width=1.1pt] (K) -- ($(K)!0.93!(M)$);
  \draw[-{Latex[length=2mm]},Blue,line width=1.1pt] (K) -- ($(K)+(-45:12)$);
  \fill[Navy] (K) circle (1.8);
  \node[Force,font=\small] at ($(K)!0.5!(M)+(-2.8,2.8)$) {$\FZ$};
  \node[Blue,font=\small,anchor=west] at ($(K)+(-45:12)+(0.4,0)$) {$\vB$};
  \node[font=\scriptsize,Muted,anchor=north] at (0,-17.2) {mit $\FZ$: Kreisbahn};
\end{tikzpicture}
\end{minipage}\hfill
\begin{minipage}[c]{0.24\linewidth}\centering
\begin{tikzpicture}[x=1mm,y=1mm]
  \coordinate (M) at (0,0);
  \coordinate (K) at (45:15);
  \draw[Muted,line width=.6pt] (M) circle (15);
  \fill[Navy] (M) circle (.8);
  \node[font=\scriptsize,anchor=north] at (0,-1.2) {Mittelpunkt};
  \draw[Blue,densely dashed,line width=.9pt] (K) -- ($(K)+(-45:20)$);
  \draw[-{Latex[length=2mm]},Blue,line width=1.1pt] (K) -- ($(K)+(-45:12)$);
  \fill[Navy] (K) circle (1.8);
  \node[Blue,font=\small,anchor=west] at ($(K)+(-45:12)+(0.4,0)$) {$\vB$};
  \node[font=\scriptsize,Muted,anchor=north] at (0,-17.2) {ohne $\FZ$: geradlinig weiter};
\end{tikzpicture}
\end{minipage}\hfill
\begin{minipage}[c]{0.48\linewidth}
\small\textcolor{Blue}{\bfseries Musterlösung Teilaufgabe 7:} Der Geschwindigkeitspfeil beginnt am Körper und liegt tangential zur Kreisbahn. Der Pfeil der Zentripetalkraft $\FZ$ beginnt ebenfalls am Körper und zeigt \textbf{nach innen} zum Kreismittelpunkt.

\vspace{1.2mm}
Fällt $\FZ$ weg, bewegt sich der Körper wegen seiner Trägheit \textbf{tangential geradlinig} weiter.
\end{minipage}

\vspace{2mm}
\begin{tcolorbox}[colback=Mint,colframe=Blue,arc=1.2mm,boxrule=.7pt,left=2.5mm,right=2.5mm,top=1.2mm,bottom=1.2mm]
\textcolor{Blue}{\bfseries Merksatz}\par\vspace{0.6mm}
Wirkt eine \textbf{Kraft} auf einen Körper, ändert sich \textbf{der Betrag der Geschwindigkeit} und/oder \textbf{die Richtung der Geschwindigkeit}.

\vspace{0.8mm}
Bei einer \textbf{Kreisbewegung} wirkt die Zentripetalkraft nach \textbf{innen}, denn es ändert sich ständig \textbf{die Richtung der Geschwindigkeit}.
\end{tcolorbox}

\vfill
{\color{Line}\rule{\linewidth}{0.5pt}}\par\vspace{1mm}
\begin{tabularx}{\linewidth}{@{}Xr@{}}
  {\small\color{Muted}Klasse 11 · Physik · Kreisbewegungen} &
  {\small\color{Muted}Sicherungsblatt · Aufgabe 3}
\end{tabularx}

\end{document}
